% !TEX root = ../../hw07.tex
% Source: author's supplied text in chat, 2026-11-06, Q6.
% The supplied proof establishes decrease but omits the requested limit.
% The added logarithmic limit and other amendments were approved, 2026-11-06.
\begin{proof}[Solution]
  To show the function
  \[
    f(x)=\left(1+\frac1x\right)^{x+1},
    \qquad x>0,
  \]
  decreases, we can analyze the derivative of its logarithm, since
  $\ln f(x)$ preserves the order and monotonicity of $f$.
  The base is positive, so $f(x)>0$. Set
  $H(x)=\ln f(x)=(x+1)\ln(1+1/x)$. Then
  \[
    \begin{aligned}
      H'(x)&=\ln\left(1+\frac1x\right)
        +(x+1)\frac1{1+\frac1x}\left(-\frac1{x^2}\right)\\
        &=\ln\left(1+\frac1x\right)-\frac{x+1}{x^2+x}\\
        &=\ln\left(1+\frac1x\right)-\frac1x.
    \end{aligned}
  \]
  Let $\varphi(u)=\ln(1+u)-u$, where $u=1/x$ and $u>0$.
  Since $\varphi$ is
  continuous on $[0,u]$ and differentiable on
  $(0,u)$, with $\varphi(0)=0$, we can apply the mean value theorem
  to find a $\xi\in(0,u)$ such that
  \[
    \begin{aligned}
      \frac{\varphi(u)-\varphi(0)}{u-0}&=\varphi'(\xi),\\
      \varphi(u)&=u\varphi'(\xi).
    \end{aligned}
  \]
  Since $\xi>0$, the derivative becomes negative,
  \[
    \varphi'(\xi)=\frac1{1+\xi}-1<0.
  \]
  A negative $\varphi'(\xi)$ implies $\varphi(u)<0$ for all $u>0$.
  Therefore, applying
  \[
    H'(x)=\varphi(u)=\varphi(1/x),\qquad
      x>0,
  \]
  we can conclude that $H'(x)<0$ and therefore $f$ is strictly
  decreasing.
  \autoref{fig:hw07:decreasing-approximants}
  shows this decrease for $x\ge1$.
  \begin{marginfigure}
    \centering
    \input{\analysisroot/figures/interval-styles.tex}
\begin{tikzpicture}[interval diagram,x=.42cm,y=1cm]
  \draw[->] (0,0)--(8.8,0) node[right=4pt] {$x$};
  \draw[->] (0,0)--(0,4.5) node[above=4pt] {$y$};
  \node[below left=5pt] at (0,0) {$0$};
  \draw[gray,dashed] (0,{exp(1)})--(8.4,{exp(1)});
  \node[left=5pt] at (0,{exp(1)}) {$e$};
  \draw (0,4)--(-.13,4) node[left=5pt] {$4$};
  \draw[->,domain=1:8.4,samples=90]
    plot (\x,{exp((\x+1)*ln(1+1/\x))});
  \foreach \n in {1,4,8} {
    \fill (\n,{exp((\n+1)*ln(1+1/\n))}) circle[radius=1.6pt];
    \draw (\n,0)--(\n,-.06) node[below=5pt] {$\n$};
  }
  \draw (.2,-.95)--(1.05,-.95);
  \node[anchor=west] at (1.35,-.95) {$f$};
  \draw[gray,dashed] (.2,-1.4)--(1.05,-1.4);
  \node[anchor=west] at (1.35,-1.4) {limiting level $e$};
\end{tikzpicture}

    \caption[Decreasing exponential approximants.]{%
      The actual graph for $x\ge1$ descends toward $e$. The limit
      calculation establishes the dashed level. Since $f$ is
      strictly decreasing on $x>0$, every value remains above $e$.}
    \label{fig:hw07:decreasing-approximants}
  \end{marginfigure}

  To find the limit, put $t=1/x$. As
  $x\to+\infty$, we have $t\to0^+$ and
  \[
    H(x)
      =\frac{\ln(1+t)}{t}+\ln(1+t).
  \]
  The quotient is of $0/0$ form, and its
  denominator derivative is $1$. L'H\^opital's rule therefore gives
  \[
    \lim_{t\to0^+}\frac{\ln(1+t)}{t}
      =\lim_{t\to0^+}\frac1{1+t}=1.
  \]
  Continuity of the logarithm also gives
  $\ln(1+t)\to0$. Thus $H(x)\to1$, and continuity of the
  exponential recovers the requested limit,
  \[
    \lim_{x\to+\infty}
      \left(1+\frac1x\right)^{x+1}=e^1=e.
  \]
\end{proof}
